\chapter{Allergy Skin Test}

\section*{What Is This Test?}
An allergy skin test identifies immediate hypersensitivity to selected allergens by observing the skin's reaction. Skin-prick testing is most common; intradermal testing may be used for selected medicines or insect venoms, and patch testing evaluates delayed contact dermatitis. A positive result shows sensitization, not necessarily that the allergen causes symptoms in daily life.

\section*{Alternative Names}
\begin{itemize}
  \item Skin-Prick Test
  \item Scratch Test
  \item Puncture Test
  \item Intradermal Allergy Test
  \item Patch Test
  \item Skin Allergy Testing
\end{itemize}
\section*{How This Test Is Performed}
Small drops of allergens are placed on the forearm or back and the skin is lightly pricked. After about 15--20 minutes, wheals and redness are measured against positive and negative controls. Patch testing uses allergen-containing adhesive chambers that remain in place for 48--96 hours and are read later. Testing should be performed by a trained clinician with treatment available for a systemic reaction.

\section*{Why It Is Ordered}
Evaluation of suspected environmental, food, venom, or selected medication allergy; allergic rhinitis or asthma; recurrent hives; eczema or contact dermatitis; and planning of allergen avoidance or immunotherapy.

\section*{Preparation, Risks, and Limitations}
Antihistamines and some other medicines can suppress the response and may need to be withheld only as directed by the clinician. Skin with extensive eczema or dermatographism may be difficult to interpret. Local itching and swelling are expected; a systemic reaction is uncommon but possible. Results must be interpreted with the symptom history and should not be used to diagnose an allergy without compatible clinical symptoms.

\section*{References}
\begin{enumerate}
  \item MedlinePlus. Allergy Skin Test.
  \item American College of Allergy, Asthma \& Immunology. Allergy testing guidance.
\end{enumerate}

\clearpage
\section*{Understanding Results and Follow-up}
The laboratory report should identify the specimen, method, units, reference interval or decision limit, and whether the result is positive, negative, abnormal, or inconclusive. Reference intervals vary with age, sex, pregnancy, specimen type, assay, and laboratory; a result outside a range is not automatically a diagnosis, and a result within a range does not always exclude disease. Discuss unexpected results with the ordering clinician, who can decide whether repeat or confirmatory testing, treatment, or referral is appropriate. Do not change medicines or delay urgent care based on an isolated result.


\section*{Regulatory Status and Authorization Date}
This chapter describes a test category, not a specific commercial product. FDA, CDSCO, EU IVDR, and MHRA approval or registration dates therefore cannot be assigned without the assay manufacturer, product name, instrument, and intended use. For laboratory-developed tests, an individual agency approval date may not exist. See the Preface for the interpretation of regulatory status.

\clearpage
